\documentclass[11pt]{article}

\usepackage[a4paper,margin=1in]{geometry}
\usepackage{amsmath,amssymb,amsthm}
\usepackage[T1]{fontenc}
\usepackage{hyperref}

\newtheorem{lemma}{Lemma}[section]
\newtheorem{theorem}[lemma]{Theorem}
\newtheorem{remark}[lemma]{Remark}

\begin{document}

\begin{center}
{\Large\bfseries PRIMES IN A LOGARITHMIC BLOCK PRODUCT}
\end{center}

\medskip

\noindent\textsc{Abstract.} For $n \ge 3$ let
\[
F(n):=\prod_{1\le i\le \lfloor \log n\rfloor}(n+i),
\]
and let $q(n,k)$ denote the least prime that does not divide $\prod_{1\le i\le k}(n+i)$. We prove that for every fixed real $C>2$ there are infinitely many integers $n$ such that every prime $p\le C\log n$ divides $F(n)$; equivalently, $q(n,\lfloor \log n\rfloor)>C\log n$ infinitely often. The construction is elementary apart from the standard estimate $\theta(x)=x+o(x)$ for the Chebyshev function. It combines a multiple of $\prod_{k<p\le 2k}p$ near $e^k$ with a weighted simultaneous approximation argument for the primes in $(2k,Ak]$, where $A>C$ is fixed. We also record the unconditional upper bound $q(n,\lfloor \log n\rfloor)\le (1+o(1))(\log n)^2$.

\section{Introduction}

For $n\ge 3$ define
\[
F(n):=\prod_{1\le i\le \lfloor \log n\rfloor}(n+i),
\]
where $\log$ denotes the natural logarithm. The problem considered here asks whether there exists some $\varepsilon>0$ such that infinitely many integers $n$ satisfy
\[
\forall p\le (2+\varepsilon)\log n,\qquad p\mid F(n),
\]
with $p$ running over the primes. Equivalently, if $q(n,k)$ denotes the least prime that does not divide $\prod_{1\le i\le k}(n+i)$, the question is whether
\[
q(n,\lfloor \log n\rfloor)>(2+\varepsilon)\log n
\]
for infinitely many $n$.

We prove a stronger statement: the constant $2+\varepsilon$ may be replaced by an arbitrary fixed constant $C>2$. The proof has two ingredients. First, if
\[
Q_k:=\prod_{k<p\le 2k}p,
\]
then a multiple of $Q_k$ automatically supplies divisibility by every prime in $(k,2k]$. Second, after multiplying by a factor of size $\exp(o(k))$, one can arrange that the resulting center lies within distance $(k-1)/2$ of a multiple of every prime in $(2k,Ak]$, where $A>C$ is fixed. Since the center is chosen near $e^k$, the corresponding block of $k$ consecutive integers is contained in the longer block defining $F(n)$, and its logarithmic scale is still $k+o(k)$.

We use the standard notation
\[
\theta(x):=\sum_{p\le x}\log p.
\]

The only asymptotic input is the prime number theorem in the form
\begin{equation}
\theta(x)=x+o(x)\qquad (x\to\infty).
\end{equation}

\bigskip
\noindent\rule{2cm}{0.4pt}

\noindent\textit{2020 Mathematics Subject Classification.} 11A41, 11B25, 11J25.

\noindent\textit{Key words and phrases.} prime divisors, short intervals, simultaneous approximation, pigeonhole principle.

\section{The lower bound construction}

We first isolate the simultaneous approximation step in a convenient weighted form.

\begin{lemma}
Let $r\ge 1$, let $\alpha_1,\dots,\alpha_r\in\mathbb{R}$, and let $M_1,\dots,M_r$ be positive integers. Put
\[
B:=\prod_{j=1}^r M_j.
\]
Then there exists an integer $t$ with $1\le t\le B$ such that
\[
\|t\alpha_j\|\le \frac{1}{M_j}\qquad (1\le j\le r),
\]
where $\|x\|$ denotes the distance from $x$ to the nearest integer.
\end{lemma}

\begin{proof}
Consider the $B+1$ points of $[0,1)^r$
\[
x_u:=\bigl(\{u\alpha_1\},\dots,\{u\alpha_r\}\bigr),\qquad 0\le u\le B,
\]
where $\{x\}$ denotes the fractional part of $x$. Partition $[0,1)^r$ into the half-open boxes
\[
\prod_{j=1}^r \left[\frac{a_j}{M_j},\frac{a_j+1}{M_j}\right),\qquad 0\le a_j\le M_j-1.
\]
There are exactly $B$ such boxes. Hence two of the points, say $x_u$ and $x_v$ with $u>v$, lie in the same box. Put $t:=u-v$. Then $1\le t\le B$, and for every $j$ the two fractional parts $\{u\alpha_j\}$ and $\{v\alpha_j\}$ differ by at most $1/M_j$ inside the unit circle. Equivalently,
\[
\|t\alpha_j\|\le \frac{1}{M_j}\qquad (1\le j\le r).
\]
\end{proof}

The main theorem is as follows.

\begin{theorem}
For every fixed $C>2$, there exist infinitely many integers $n$ such that every prime
\[
p\le C\log n
\]
divides $F(n)$. Equivalently,
\[
q(n,\lfloor \log n\rfloor)>C\log n
\]
for infinitely many $n$.
\end{theorem}

\begin{proof}
Fix $C>2$, and choose a real number $A>C$. Let $k$ be a large odd integer, and write
\[
h:=\frac{k-1}{2}.
\]
Define
\[
Q_k:=\prod_{k<p\le 2k}p.
\]
By (1),
\begin{equation}
\log Q_k=\theta(2k)-\theta(k)=k+o(k).
\end{equation}
Let $M_k$ be the least positive multiple of $Q_k$ that is not smaller than $e^{k+2}$. Then
\[
Q_k\mid M_k,\qquad M_k\ge e^{k+2},\qquad M_k<e^{k+2}+Q_k.
\]
In particular, (2) implies
\begin{equation}
\log M_k=k+o(k).
\end{equation}

Let $p_1,\dots,p_r$ be the primes in the interval $(2k,Ak]$. For each $j$ put
\[
N_j:=\left\lfloor \frac{p_j}{h}\right\rfloor,\qquad B_k:=\prod_{j=1}^r N_j,
\]
with the convention $B_k=1$ if $r=0$. Since $p_j\le Ak$ and $h=(k-1)/2$, we have
\[
N_j\le \frac{p_j}{h}+1\le \frac{2Ak}{k-1}+1\le 2A+2
\]
for all sufficiently large $k$. Also, by the prime number theorem,
\[
r=\pi(Ak)-\pi(2k)=O\!\left(\frac{k}{\log k}\right).
\]
Therefore
\begin{equation}
\log B_k=O(r)=O\!\left(\frac{k}{\log k}\right)=o(k).
\end{equation}

If $r\ge 1$, apply Lemma 2.1 with
\[
\alpha_j=\frac{M_k}{p_j},\qquad M_j=N_j.
\]
We obtain an integer $t_k$ with
\[
1\le t_k\le B_k
\]
and
\[
\left\|\frac{t_kM_k}{p_j}\right\|\le \frac{1}{N_j}\le \frac{h}{p_j}\qquad (1\le j\le r).
\]
Thus for each $j$ there exists an integer $m_j$ such that
\begin{equation}
|t_kM_k-m_jp_j|\le h.
\end{equation}
If $r=0$, set simply $t_k:=1$. In either case,
\begin{equation}
1\le t_k\le B_k,\qquad \log t_k=o(k)
\end{equation}
by (4).

Now put
\[
m_k:=M_kt_k,\qquad n_k:=m_k-h-1.
\]
Since $k=2h+1$, the block
\[
n_k+1,\dots,n_k+k
\]
is exactly the symmetric interval
\begin{equation}
[m_k-h,m_k+h].
\end{equation}
We claim that every prime $p\le Ak$ divides
\[
\prod_{i=1}^k(n_k+i).
\]
There are three cases.

\textit{Case 1:} $p\le k$. Any interval of $k$ consecutive integers contains a multiple of $p$, because among any $p$ consecutive integers there is exactly one multiple of $p$, and here $k\ge p$. Hence some $n_k+i$ with $1\le i\le k$ is divisible by $p$.

\textit{Case 2:} $k<p\le 2k$. By definition, $p$ divides $Q_k$, hence also $M_k$ and therefore $m_k=M_kt_k$. Since $m_k\in [m_k-h,m_k+h]$, the integer $m_k$ itself is one of the $k$ terms in (7). Thus $p$ divides one of the factors $n_k+i$.

\textit{Case 3:} $2k<p\le Ak$. Then $p=p_j$ for some $j$. By (5), there is an integer $m_jp_j$ with
\[
|m_k-m_jp_j|\le h.
\]
Hence $m_jp_j\in [m_k-h,m_k+h]$, so again $p$ divides one of the integers $n_k+1,\dots,n_k+k$.

We have proved that every prime $p\le Ak$ divides the shorter product $\prod_{i=1}^k(n_k+i)$. We now compare $k$ with $\log n_k$. By (3) and (6),
\[
\log m_k=\log M_k+\log t_k=k+o(k).
\]
Since $m_k\ge M_k\ge e^{k+2}$ and $n_k=m_k-h-1$, we have
\[
0\le \frac{m_k-n_k}{m_k}=\frac{h+1}{m_k}\le \frac{k+1}{2e^{k+2}},
\]
so in particular
\begin{equation}
\log n_k=\log m_k+o(1)=k+o(k).
\end{equation}
The lower bound $n_k\ge e^{k+2}-(k+1)/2-1$ shows that $\log n_k\ge k+1$ for all sufficiently large $k$, and hence
\begin{equation}
\lfloor \log n_k\rfloor\ge k
\end{equation}
for all sufficiently large $k$.

Finally, since $A>C$ and (8) gives $\log n_k=k+o(k)$, we obtain
\[
C\log n_k<Ak
\]
for all sufficiently large $k$. Therefore every prime $p\le C\log n_k$ is at most $Ak$, hence divides the shorter product $\prod_{i=1}^k(n_k+i)$, and a fortiori divides
\[
F(n_k)=\prod_{1\le i\le \lfloor \log n_k\rfloor}(n_k+i)
\]
by (9). Since $n_k\to\infty$ as $k\to\infty$ through odd integers, this gives infinitely many such $n$.
\end{proof}

\begin{remark}
The proof yields more than the original formulation. Since $C>2$ is arbitrary, for every fixed $C>2$ there are infinitely many integers $n$ with
\[
q(n,\lfloor \log n\rfloor)>C\log n.
\]
Thus, the barrier at the constant $2$ can be exceeded by any prescribed amount.
\end{remark}

\begin{thebibliography}{9}

\bibitem{erdos}
P. Erd\H{o}s, \textit{Some unconventional problems in number theory}, Acta Math. Acad. Sci. Hungar. \textbf{33} (1979), 73--82.

\bibitem{bloom}
T. F. Bloom, \textit{Erd\H{o}s Problem \#457}, \url{https://www.erdosproblems.com/457} (accessed 2026-03-02).

\end{thebibliography}

\end{document}